\documentclass[border=10pt]{standalone}
\usepackage{fontspec}
\setmainfont{Noto Serif}[Ligatures=TeXOff]
\XeTeXgenerateactualtext=1
\newfontfamily\bengalifont{Noto Serif Bengali}[Script=Bengali, NFSSFamily=bnfam]
\usepackage[Latin,Bengali]{ucharclasses}
% group-free font switching (as in header-common.tex): robust inside braces, cells and quotes
\makeatletter
\def\BN@saved{\rmdefault}
\def\BN@bn{bnfam}
\newcommand\BNin{\edef\BN@tmp{\f@family}\ifx\BN@tmp\BN@bn\else\let\BN@saved\BN@tmp\fi\fontfamily{bnfam}\selectfont}
\newcommand\BNout{\fontfamily{\BN@saved}\selectfont}
\makeatother
\setTransitionsFor{Bengali}{\BNin}{\BNout}
% Typographic quotes and dashes are not in a ucharclasses class (class 0), so a Bangla word right
% after “ or ‘ would not trigger the Bangla font. Treat them as Latin so that every switch fires.
\XeTeXcharclass"2018=\XeTeXcharclass`A
\XeTeXcharclass"2019=\XeTeXcharclass`A
\XeTeXcharclass"201C=\XeTeXcharclass`A
\XeTeXcharclass"201D=\XeTeXcharclass`A
\XeTeXcharclass"2013=\XeTeXcharclass`A
\XeTeXcharclass"2014=\XeTeXcharclass`A
\XeTeXcharclass"2026=\XeTeXcharclass`A
\usepackage{newunicodechar}
\newunicodechar{।}{{\bengalifont\char"0964}}
\newunicodechar{॥}{{\bengalifont\char"0965}}

\usepackage[table]{xcolor}
\usepackage{booktabs,tabularx,array}
\definecolor{bapilrule}{HTML}{1F2937}
\definecolor{bapilstripe}{HTML}{F3F4F6}
\renewcommand{\arraystretch}{1.28}
\setlength{\tabcolsep}{5pt}

\begin{document}
\begin{minipage}{17cm}
{\large\bfseries সারণি ৬৩।}\ {\large মূল্যায়ন প্রটোকল ও রিলিজ-গেটের নিয়ম}\par\vspace{7pt}
\small
\rowcolors{2}{bapilstripe}{white}
\begin{tabularx}{\linewidth}{@{}>{\hsize=0.714\hsize\raggedright\arraybackslash}X>{\hsize=1.238\hsize\raggedright\arraybackslash}X>{\hsize=1.048\hsize\raggedright\arraybackslash}X>{\hsize=0.952\hsize\raggedright\arraybackslash}X>{\hsize=1.048\hsize\raggedright\arraybackslash}X@{}}
\arrayrulecolor{bapilrule}\toprule
\rowcolor{white}\textbf{পরিমাপ} & \textbf{সংজ্ঞা} & \textbf{পরীক্ষা-সেট ও নমুনা} & \textbf{পদ্ধতি ও দায়িত্ব} & \textbf{সিদ্ধান্ত-নিয়ম} \\
\midrule
Recall@10, MRR@10, nDCG@10 & প্রতিটি উত্তরযোগ্য প্রশ্নে কোনো সঠিক অনুচ্ছেদ প্রথম ১০টিতে থাকলে সফল; সফলতার গড়ই Recall@10। MRR ও nDCG সারণি ২৮ অনুযায়ী & ব্যক্তিগত পরীক্ষা-অংশ: v0-তে কমপক্ষে ৩০০টি উত্তরযোগ্য প্রশ্ন (প্রতি ভাষায় প্রায় ১৫০টি) এবং v1-এ ১,২০০টি & মূল্যায়ন দলের স্বয়ংক্রিয় স্কোরিং; স্ক্রিপ্ট, ফলাফল, কনফিগারেশন ও কর্পাস সংস্করণ সংরক্ষিত & বিন্দু-প্রাক্কলন সীমা অর্জন করলে এবং ৯৫\% ব্যবধান ধনাত্মক-ঋণাত্মক ৫ পয়েন্টের বেশি চওড়া না হলে সীমা অর্জিত (n = ৩০০-এ প্রায় ৪ পয়েন্ট); অন্যথায় সিদ্ধান্তের আগে আরও প্রশ্ন টীকাকরণ \\
উদ্ধৃতির নির্ভুলতা & নমুনাভুক্ত উত্তর-বাক্যের মধ্যে যেগুলোর উদ্ধৃত অনুচ্ছেদ বাক্যটিকে পুরোপুরি সমর্থন করে তার অনুপাত (সমর্থিত, আংশিক সমর্থিত বা অসমর্থিত; শুধু সমর্থিত গণ্য) & পরীক্ষা-উত্তর থেকে দৈবভাবে নেওয়া ৪০০টি উত্তর-বাক্য, প্রতি ভাষায় কমপক্ষে ২০০টি & ব্যবস্থার সংস্করণ না জানা দুজন প্রশিক্ষিত টীকাকার; তৃতীয়জন নিষ্পত্তি করেন; কোহেনের কাপ্পা প্রতিবেদিত (লক্ষ্য ০.৭০ বা তার বেশি) & বিন্দু-প্রাক্কলন কমপক্ষে ০.৯৫ (n = ৪০০-এ ব্যবধান প্রায় ধনাত্মক-ঋণাত্মক ২ পয়েন্ট) \\
অসমর্থিত দাবির হার & নমুনাভুক্ত বাক্যের মধ্যে যেগুলো কোনো উদ্ধারকৃত অনুচ্ছেদ দ্বারা সমর্থিত নয় তার অনুপাত & উদ্ধৃতির নির্ভুলতার একই নমুনা & একই টীকাকার & বিন্দু-প্রাক্কলন সর্বোচ্চ ২\% \\
সঠিক বিরতি ও উত্তর-আওতা & আওতাবহির্ভূত প্রশ্নে বিরতি; উত্তরযোগ্য প্রশ্নে উত্তর & v0-তে ১৫০টি আওতাবহির্ভূত পরীক্ষা-প্রশ্ন (রচিত ২০০টি: উন্নয়নে ৫০, পরীক্ষায় ১৫০); সব উত্তরযোগ্য পরীক্ষা-প্রশ্ন & স্বয়ংক্রিয় & বিরতি কমপক্ষে ০.৯০ (ব্যবধান প্রায় ধনাত্মক-ঋণাত্মক ৫ পয়েন্ট); আওতা কমপক্ষে ০.৮০ \\
ভাষাগত সমতা & ইংরেজির Recall@10 বিয়োগ বাংলার Recall@10, বুটস্ট্র্যাপ ব্যবধানসহ & ভাষাভিত্তিক পরীক্ষা-উপসেট & স্বয়ংক্রিয় & v0-তে প্রতিবেদিত, কারণ ব্যবধান সিদ্ধান্তের জন্য বেশি চওড়া; v1 থেকে রিলিজের শর্ত \\
স্থিতিশীলতা & ভিন্ন বানান ও রোমান হরফের প্রশ্নে, OCR-নির্ভর নথিতে উত্তরযুক্ত প্রশ্নে এবং মিশ্র-ভাষার প্রশ্নে Recall@10 & যথাক্রমে কমপক্ষে ১০০, ১০০ ও ৫০টি পরীক্ষা-প্রশ্ন & স্বয়ংক্রিয় & বাংলা হরফে ভিন্ন বানানে v0-তে কমপক্ষে ০.৭৫; রোমান হরফ, OCR-নির্ভর ও মিশ্র-ভাষার ফলাফল v0-তে গেট ছাড়া প্রতিবেদিত; প্রতিবর্ণীকরণ চালু হওয়ার পর v1 থেকে রোমান হরফ ও মিশ্র-ভাষার প্রশ্নে গেট (সারণি ৩৭) \\
প্রম্পট-ইনজেকশন ও নিষিদ্ধ ফলাফল & রেড-টিমের প্রচেষ্টার (নথি বা প্রশ্নে লুকানো নির্দেশ) মধ্যে যেগুলো আচরণ বদলায় তার অনুপাত; নিষিদ্ধ ফলাফলের অনুরোধ & কমপক্ষে ১০০টি প্রচেষ্টা, প্রতি রিলিজে হালনাগাদ & নিরাপত্তা প্রধান ও একটি বহিঃস্থ রেড-টিম & সর্বোচ্চ ৫\%; অ-উন্মুক্ত ডেটার যেকোনো ফাঁস অনুত্তীর্ণ \\
গোপনীয়তা & লগ, ফলাফল বা মূল্যায়ন-সেটে P2 বা P3 ডেটার উপস্থিতি; প্রবেশাধিকার-নিয়ন্ত্রণ পরীক্ষা & লগের নমুনা ও স্বয়ংক্রিয় স্ক্যান & শাসন বিশ্লেষক; বার্ষিক স্বাধীন নিরীক্ষা & শূন্য অনুসন্ধান-ফল \\
লেটেন্সি ও প্রাপ্যতা & p95 প্রতিক্রিয়া-সময়; কর্মঘণ্টার যে অংশে সেবা চালু & G2-তে ৩০ জন একযোগে ব্যবহারকারীসহ চাপ-পরীক্ষা; G3 চলাকালে পর্যবেক্ষণ & প্ল্যাটফর্ম প্রধান & অনুসন্ধানে ৩ সেকেন্ড ও উত্তরে ১০ সেকেন্ড; প্রথম বছরে প্রাপ্যতা প্রতিবেদিত, দ্বিতীয় বছর থেকে কমপক্ষে ৯৯\% (NFR-04) \\
উপযোগিতা & ৫-পয়েন্ট স্কেলে গড় মূল্যায়ন & কমপক্ষে ৩০ জন পাইলট ব্যবহারকারী; বাংলা সম্পাদকদের মূল্যায়িত কমপক্ষে ১০০টি উত্তর & প্রোডাক্ট ওনার ও মূল্যায়ন দল & কমপক্ষে ৪.০ \\
P2 নিষ্কাশনের নির্ভুলতা & নিষ্কাশিত সারণি-ঘরের মধ্যে যেগুলো মূল নথির সঙ্গে মেলে তার অনুপাত & বছর ও মন্ত্রণালয়জুড়ে স্তরভিত্তিক দৈবচয়িত কমপক্ষে ১,০০০টি ঘর & দুজন স্বাধীন যাচাইকারী & কমপক্ষে ৯৯\% (ব্যবধান প্রায় ধনাত্মক-ঋণাত্মক ০.৬ পয়েন্ট) \\
প্রমাণ-প্রাপ্তির সময় & সমতুল্য কাজে সরঞ্জামসহ ও সরঞ্জাম ছাড়া প্রামাণ্য প্রমাণ খুঁজে উদ্ধৃত করার মধ্যক সময় & কমপক্ষে ২০ জন বিশ্লেষক, প্রত্যেকে ৬টি কাজ, ক্রসওভার নকশা & মূল্যায়ন দল; মাস ১৮-এ স্বাধীন মূল্যায়নকারী & প্রভাব ও ব্যবধান প্রতিবেদিত; G4-কে তথ্য দেয়, G3-তে উত্তীর্ণ-নম্বর নেই \\
পুনরুৎপাদনযোগ্যতা & ব্যক্তিগত অংশে সংরক্ষিত ফলাফলের পুনঃস্কোরিং & গেট-সিদ্ধান্তে ব্যবহৃত প্রতিটি রিলিজ & G4-এর আগে স্বাধীন মূল্যায়নকারী & পুনঃস্কোরকৃত মাপকাঠি প্রতিবেদিত মানের ১ পয়েন্টের মধ্যে \\
\bottomrule
\end{tabularx}
\par\vspace{5pt}{\footnotesize\itshape ব্যবধানগুলো ৯৫\% উইলসন বা বুটস্ট্র্যাপ ব্যবধান। BanglaPolicyQA v0-তে ৫০০টি উত্তরযোগ্য প্রশ্ন (৪০\% উন্মুক্ত উন্নয়ন-অংশ, ৬০\% ব্যক্তিগত পরীক্ষা-অংশ) ও ২০০টি আওতাবহির্ভূত প্রশ্ন আছে। নমুনার আকার ন্যূনতম; T2 স্তরের ব্যবহারে ERRC ও TAB আরও বেশি চাইতে পারে।}
\par\vspace{3pt}{\scriptsize\color{gray} উৎস: BAPIL টেকনিক্যাল ব্লুপ্রিন্ট, সংস্করণ ১.০ (২০২৬)। CC BY 4.0।}
\end{minipage}
\end{document}
